% exp01_smoke 实验结果表（由 tables.py 从 metrics.json / solution.json 生成）

% 编译：latexmk -xelatex exp01_tables.tex（ctex 需 xelatex；run.py 自动编译）

\documentclass[a4paper,11pt]{article}

\usepackage[margin=2.2cm]{geometry}

\usepackage{booktabs}

\usepackage{ctex} % 中文；推荐 xelatex 编译（无 ctex 宏包时改英文文档类）

\begin{document}

\title{exp01\_smoke 实验结果记录}

\maketitle

\begin{table}[htbp]\centering
\caption{exp01\_smoke 运行配置总览（seed=42）}
\label{tab:exp01-config}
\begin{tabular}{ll}
\toprule
配置项 & 取值 \\
\midrule
随机种子 & [42] \\
滚动周期数 & 3 \\
机巢数 / 单巢容量 & 8 / 4 \\
机队规模 $K_{\mathrm{avail}}$ & 10 \\
无人机速度 / 满电航程 & 15.0 m/s / 80.0 km \\
任务总数（首批/后续到达） & 12 + 8 \\
DMDE 种群 / 迭代代数 & 50 / 100 \\
目标标量化 & 字典序（big-M，收益优先） \\
交付时间窗约束 & 开启 \\
DEM 区域 / 重采样步长 & N29E091 / 200.0 m \\
\bottomrule
\end{tabular}
\end{table}


\begin{table}[htbp]\centering
\caption{逐周期任务漏斗与任务命运（闸门拦截 = 交付被判违约、解绑回池的任务次）}
\label{tab:exp01-funnel}
\begin{tabular}{rcccccccccc}
\toprule
周期 & 新到 & 滞留 & 释放 & 过预筛 & 选中 & 执行 & 延期 & 完成 & 拦截 \\
\midrule
1 & 12 & 0 & 0 & 12 & 12 & 12 & 0 & 11 & 1 \\
2 & 4 & 0 & 1 & 5 & 5 & 5 & 0 & 5 & 0 \\
3 & 4 & 0 & 0 & 4 & 4 & 4 & 0 & 4 & 0 \\
\bottomrule
\end{tabular}
\end{table}


\begin{table}[htbp]\centering
\caption{逐周期目标、约束可行性与求解开销（$F_1^{\mathrm{plan}}$ 为规划口径收益，$F_1^{\mathrm{comp}}$ 为闸门后实际完成收益；$C^{3D}$ 为三维路径总代价，单位 km）}
\label{tab:exp01-objective}
\begin{tabular}{rccccccc}
\toprule
周期 & $F_1^{\mathrm{plan}}$ & $F_1^{\mathrm{comp}}$ & $C^{3D}$ (km) & 惩罚 & 违约数 & 可行 & 求解 (s) \\
\midrule
1 & 60.50 & 57.79 & 485.3 & 944.80 & 1 & 否 & 79.7 \\
2 & 31.17 & 31.17 & 207.8 & 0.00 & 0 & 是 & 3.3 \\
3 & 27.95 & 27.95 & 157.9 & 0.00 & 0 & 是 & 1.9 \\
\bottomrule
\end{tabular}
\end{table}


\begin{table}[htbp]\centering
\caption{时间窗违约与出动明细（按任务列出全部曾出动任务；重复出动 = 同一任务多周期再次派出）}
\label{tab:exp01-tw}
\begin{tabular}{rcccl}
\toprule
任务 & $t_{\mathrm{latest}}$ (min) & 出动次数 & 闸门拦截 & 结局 \\
\midrule
T0 & 50.8 & 1 & 0 & 完成 \\
T1 & 66.4 & 1 & 0 & 完成 \\
T2 & 55.2 & 1 & 0 & 完成 \\
T3 & 75.1 & 1 & 0 & 完成 \\
T4 & 64.7 & 1 & 0 & 完成 \\
T5 & 82.5 & 1 & 0 & 完成 \\
T6 & 76.5 & 1 & 0 & 完成 \\
T7 & 59.1 & 1 & 0 & 完成 \\
T8 & 82.5 & 1 & 0 & 完成 \\
T9 & 81.2 & 1 & 0 & 完成 \\
T10 & 62.4 & 1 & 0 & 完成 \\
T11 & 58.0 & 2 & 1 & 完成 \\
T12 & 67.5 & 1 & 0 & 完成 \\
T13 & 67.5 & 1 & 0 & 完成 \\
T14 & 74.0 & 1 & 0 & 完成 \\
T15 & 55.2 & 1 & 0 & 完成 \\
T16 & 107.3 & 1 & 0 & 完成 \\
T17 & 106.7 & 1 & 0 & 完成 \\
T18 & 97.2 & 1 & 0 & 完成 \\
T19 & 65.9 & 1 & 0 & 完成 \\
\bottomrule
\end{tabular}
\end{table}


\begin{table}[htbp]\centering
\caption{结果总览与收益核算对账}
\label{tab:exp01-outcome}
\begin{tabular}{ll}
\toprule
指标 & 数值 \\
\midrule
到达任务总数 $|\mathcal{T}|$ & 20 \\
完成任务数 & 20（唯一 20） \\
完成率 & 100.0\% \\
周期解可行率 & 66.7\%（2/3；其中 1 个周期含时间窗违约） \\
规划口径收益 $\sum F_1^{\mathrm{plan}}$ & 119.61 \\
完成口径收益 $\sum F_1^{\mathrm{comp}}$ & 116.91 \\
核算膨胀率 & 2.3\%（闸门拦截 + 重复出动重复计数） \\
总三维路径代价 & 851.0 km \\
平均等待时间 & 10.5 min \\
总出动任务次 / 重复出动 & 21 / 1 \\
终点机巢使用 & 7 / 8 个 \\
总求解时间 & 84.9 s \\
\bottomrule
\end{tabular}

\vspace{2pt}\parbox{0.92\linewidth}{\footnotesize \textbf{口径注}：$F_1^{\mathrm{plan}}$ 为求解期目标值（含被完成闸门拦截的交付、以及重复出动任务的重复计数）；$F_1^{\mathrm{comp}}$ 按最终完成的唯一任务计。论文报完成率与收益时应使用完成口径，规划口径仅用于说明优化过程。}
\end{table}


\end{document}